\documentclass[12pt]{article}

\usepackage{amsmath}
\usepackage{amssymb}
\usepackage{graphicx}
\graphicspath{{figures/}{figures/xi_estimators/}}
\usepackage{hyperref}
\usepackage[margin=1in]{geometry}
\usepackage{xcolor}
\usepackage{tcolorbox}
\tcbuselibrary{breakable,skins}
\usepackage{natbib}
\bibliographystyle{aasjournal}

% Notation follows cosmo_stat.tex (bold vectors, \avg, \hat\delta_\bk,
% P(\bk) = \avg{|\hat\delta_\bk|^2}/V) and the Fourier sign convention of
% cosmo_ic.tex App. A (e^{-i k.x} forward).
\newcommand{\bk}{\mathbf{k}}
\newcommand{\bx}{\mathbf{x}}
\newcommand{\by}{\mathbf{y}}
\newcommand{\br}{\mathbf{r}}
\newcommand{\bd}{\mathbf{d}}
\newcommand{\bs}{\mathbf{s}}
\newcommand{\avg}[1]{\left\langle #1 \right\rangle}
\newcommand{\code}[1]{\texttt{#1}}

% Mask (selection function) of the measurement region.
\newcommand{\W}{W}
% Correlation (star) product, as used by Slepian & Eisenstein.
\newcommand{\corr}{\star}

\title{Estimating the Correlation Function\\
       on a Grid}
\author{}
\date{\today}

\begin{document}

\maketitle

\begin{tcolorbox}[colback=yellow!10!white, colframe=orange!75!black,
                  title={\bfseries Disclaimer}]
These notes were compiled by an AI assistant (Claude, Anthropic) and have
\textbf{not been reviewed by a human domain expert}.  Cross-check key
equations and conclusions before relying on them.
\end{tcolorbox}

\tableofcontents

\section{Two ways to compute $\xi$ on a grid}\label{sec:question}

You have a box of particles, or a patch cut out of one, and you want the
correlation function of the density field inside it.  One program walks over
every pair of grid cells at a given separation, multiplies the two density
contrasts, and averages.  Another Fourier transforms the field, squares the
modulus, and transforms back.  Run both on the same data and they print the
same number, digit after digit, until the last two or three.

The first program is doing what the definition of $\xi$ says, so its answer
needs no explanation.  The second one does.  A Fourier transform mixes every
cell into every coefficient, so a value at one separation appears to depend on
data from the far side of the box.  Squaring the modulus discards the phases,
and the phases are most of what tells one field from another.  An FFT is also
a numerical algorithm with its own rounding behaviour, so even if it computed
the right quantity one would expect an approximation to it rather than the
same number.

None of those objections holds up, and the reason is worth working through
slowly.  The Fourier method does not approximate the lag sum.  It
is the same finite sum with its terms added in a different order, so in exact
arithmetic the two are one expression.  That leaves only the last two or three
digits to account for, which \S\ref{sec:roundoff} does.

The two programs are real ones.  They are
\texttt{measure\_xi\_mpi.c}, which evaluates the direct sum, and
\texttt{measure\_xi\_fft\_mpi.c}, which evaluates the padded FFT.  Both
sit in the KIAS \texttt{HR-ICs-pipeline} repository and are documented as
agreeing bin for bin.  Everything below stands on its own without them.

\citet{Slepian2016} give the compact statement of all this for galaxy surveys:
the Landy--Szalay estimator is itself a binned convolution, so its numerator
and its denominator can both be evaluated with fast Fourier transforms in
$\mathcal{O}(N_g \log N_g)$ operations.  Their \S2 is the reference for
\S\ref{sec:ls}, \S\ref{sec:pad-howmuch} and \S\ref{sec:gives-up} below, and
they go on to the anisotropic two-point function and the three-point function,
which this note does not.

\section{The gridded field, the mask, and the estimator}\label{sec:setup}

\subsection{The gridded density contrast}

Start from particles inside a rectangular measurement region.  Lay a uniform
grid of cells over that region and assign each particle's mass to nearby
cells using some mass-assignment scheme; cloud-in-cell (CIC) is the usual
choice, and its details do not matter until \S\ref{sec:gives-up}.  The result
is a mass per cell, which divided by the cell volume is a density $\rho_x$,
where $x$ now labels a cell rather than a continuous position.

Define the density contrast in the usual way,
\begin{equation}\label{eq:delta-def}
  \delta_x \;=\; \frac{\rho_x}{\bar\rho} - 1 ,
\end{equation}
where $\bar\rho$ is the mean density \emph{measured over the region itself}.
The choice of $\bar\rho$ matters and will come back in
\S\ref{sec:integral-constraint}; for now note only that it forces
$\sum_x \delta_x = 0$ when the sum runs over the region.

Both estimators discussed below start from the same array $\delta_x$.  Nothing
in \S\S\ref{sec:ls}--\ref{sec:roundoff} depends on how that array was
produced.

\subsection{The mask}

We need a way to say that one cell is inside the measurement region and
another is not.  Introduce the \emph{mask}, sometimes called the selection
function,
\begin{equation}\label{eq:mask}
  \W_x \;=\;
  \begin{cases}
    1 & \text{cell $x$ is inside the region},\\
    0 & \text{otherwise}.
  \end{cases}
\end{equation}
The name ``mask'' is borrowed from survey analysis, where $\W$ encodes which
patches of sky were actually observed.  Here the region is a rectangular box
cut out of a larger simulation, so $\W$ is the indicator function of a
rectangle, but it is worth keeping the general notation because
\S\ref{sec:denominator} turns on the fact that our $\W$ is rectangular.

Zero-padding a field outside the region --- literally storing zeros there ---
is the same thing as forming the product $\W_x \delta_x$.  That observation is
the entire content of the padding trick, and it is worth stating early.

\subsection{The estimator}

The word \emph{estimator} is used throughout what follows, so it is worth
saying plainly what it means.  An estimator is a formula that turns the data
you actually have into a number meant to approximate a quantity you cannot
observe.  Here the data is the gridded density field in one box.  The
quantity is $\xi(r)$, which by its definition is an average over an ensemble
of universes, and there is only one universe to look at.  The formula below
is therefore not $\xi$ itself but a stand-in for it, and the hat in
$\hat\xi$ is there to keep the two apart.  Section~\ref{sec:estimator-vs-ensemble}
returns to the distinction, and \code{cosmo\_stat.pdf} treats it at length.

Two different formulas can estimate the same quantity, and two different
programs can evaluate the same formula.  Both distinctions matter below.  The
whole of this note is about one formula evaluated two ways; how well that
formula recovers the true $\xi$ is a separate question, and not one taken up
here.

We want, for each separation $d$, the average of $\delta_x \delta_{x+d}$ over
all cell pairs that are both inside the region:
\begin{equation}\label{eq:estimator}
  \hat\xi(d) \;=\; \frac{s(d)}{n(d)},
  \qquad
  s(d) = \sum_x \W_x \W_{x+d}\, \delta_x \delta_{x+d},
  \qquad
  n(d) = \sum_x \W_x \W_{x+d} .
\end{equation}
The numerator $s(d)$ is the sum of products over valid pairs.  The denominator
$n(d)$ is the number of valid pairs, which is just the numerator with the
field replaced by unity.  Both sums run over all cells of whatever array we
are storing.

Two names are worth attaching here.  The quantity $s(d)$ is the
\emph{autocorrelation} of the masked field, and $n(d)$ is the
autocorrelation of the mask.  The word autocorrelation means precisely ``a
field correlated against a shifted copy of itself,'' which is what both sums
are.

\subsection{Estimator versus ensemble quantity}\label{sec:estimator-vs-ensemble}

It is worth being clear about what $\hat\xi$ is an estimate \emph{of}, since
the distinction is the subject of the companion note \code{cosmo\_stat.pdf}
and the hat is doing real work.

The quantity we would like is the ensemble correlation function, defined
there as
\begin{equation}\label{eq:xi-ensemble}
  \xi(\br) \;\equiv\; \avg{\delta(\bx)\,\delta(\bx+\br)} ,
\end{equation}
with $\avg{\cdot}$ an average over realisations of the universe, and the
companion power spectrum
\begin{equation}\label{eq:pk-ensemble}
  P(\bk) \;=\; \frac{\avg{|\hat\delta_\bk|^2}}{V} ,
\end{equation}
where $\hat\delta_\bk = \int_V \mathrm{d}^3x\, \delta(\bx)\,
e^{-i\bk\cdot\bx}$ is the Fourier coefficient in the volume $V$, with the
sign convention of \code{cosmo\_ic.pdf} App.~A.

We have one realisation, not an ensemble, so Eq.~\eqref{eq:estimator}
replaces the ensemble average by an average over positions $x$ within the
region.  That substitution is justified by ergodicity and is exact only in the
limit of an infinite volume; at finite volume it leaves the cosmic variance
and the integral constraint discussed in
\S\ref{sec:integral-constraint}.  None of that is at issue in this note.  Both
methods estimate the \emph{same} $\hat\xi$ from the \emph{same} realisation,
so every statement below is about two evaluations of one estimator, not about
how well that estimator recovers Eq.~\eqref{eq:xi-ensemble}.

Everything that follows is about computing $s(d)$ two different ways and
about recognising $n(d)$ in closed form.

\section{Why this is the Landy--Szalay estimator}\label{sec:ls}

Equation~\eqref{eq:estimator} looks like a naive pair average, not like the
Landy--Szalay estimator \citep{LandySzalay1993} that one normally uses for a
region with edges.  In fact it is Landy--Szalay, and seeing this explains why
no random catalogue appears anywhere in the code.

\subsection{Uniform selection}\label{sec:ls-uniform}

Landy--Szalay is usually written in terms of pair counts,
\begin{equation}\label{eq:ls}
  \hat\xi_{\rm LS}(d) \;=\; \frac{DD(d) - 2\,DR(d) + RR(d)}{RR(d)} ,
\end{equation}
where $D$ is the data catalogue, $R$ is a synthetic random catalogue filling
the same region with no clustering, and $DD$, $DR$, $RR$ count pairs at
separation $d$ within and between the two catalogues.  The role of $R$ is to
measure the geometry: $RR$ tells you how many pairs at separation $d$ the
region could possibly supply, so that the shape of the region does not get
mistaken for clustering.

Now put both catalogues on our grid, and let the random catalogue become
infinitely dense while its total mass is held fixed.  In that limit $R$ stops
fluctuating and becomes exactly its own expectation value, which is a
constant density inside the region and zero outside:
\begin{equation}\label{eq:R-limit}
  R_x \;\longrightarrow\; \bar\rho\, \W_x .
\end{equation}
This is the step that removes the random catalogue from the problem.  A finite
random catalogue would carry Poisson noise that we would have to average
down by making it large; an infinitely dense one has no noise, and we can
write its contribution analytically.

Write the data field as $D_x = \bar\rho\,\W_x (1 + \delta_x)$, which is
just Eq.~\eqref{eq:delta-def} rearranged.  The three pair counts at
separation $d$ then become, using the symmetric definition of $DR$,
\begin{align}
  DD(d) &= \bar\rho^2 \sum_x \W_x \W_{x+d} (1+\delta_x)(1+\delta_{x+d}),
          \label{eq:DD}\\
  DR(d) &= \tfrac{1}{2}\bar\rho^2 \sum_x \W_x \W_{x+d}
           \big[(1+\delta_x) + (1+\delta_{x+d})\big],
          \label{eq:DR}\\
  RR(d) &= \bar\rho^2 \sum_x \W_x \W_{x+d} .
          \label{eq:RR}
\end{align}
Substituting these into Eq.~\eqref{eq:ls}, the prefactors $\bar\rho^2$
cancel, and the bracket in the numerator collapses because
\begin{equation}\label{eq:ls-algebra}
  (1+\delta_x)(1+\delta_{x+d})
  \;-\; (1+\delta_x) \;-\; (1+\delta_{x+d}) \;+\; 1
  \;=\; \delta_x \delta_{x+d} .
\end{equation}
Every term linear in $\delta$ cancels against the $DR$ subtraction, and the
constants cancel against $RR$.  What is left is
\begin{equation}\label{eq:ls-equals-ours}
  \hat\xi_{\rm LS}(d)
  \;=\; \frac{\sum_x \W_x \W_{x+d}\, \delta_x \delta_{x+d}}
             {\sum_x \W_x \W_{x+d}}
  \;=\; \frac{s(d)}{n(d)},
\end{equation}
which is exactly Eq.~\eqref{eq:estimator}.

So a grid of $\delta$ with a lag-by-lag pair-count denominator \emph{is} the
Landy--Szalay estimator, not an approximation to it.  The $RR$ term has not
been dropped; it has been evaluated analytically and appears as $n(d)$.  This
is why the code never draws a random catalogue: drawing one would only
re-introduce, with Poisson noise, a quantity we already know exactly.

\subsection{The general form: data minus randoms}\label{sec:ls-general}

The derivation above used Eq.~\eqref{eq:R-limit}, which assumes the selection
is \emph{uniform} inside the region: every cell inside the mask has the same
expected density.  That is true for a rectangular cut through a simulation
box, which is our case, and false for a real galaxy survey, where the
selection function varies with position and depth.

It is worth knowing the general form, both because it shows what our special
case is a special case \emph{of}, and because it is the version that
generalises to the Fourier method without any assumption about the geometry.
Following \citet{Slepian2016}, define the \emph{data-minus-randoms} field
\begin{equation}\label{eq:N-field}
  N \;\equiv\; D - R ,
\end{equation}
where the random catalogue is a Monte Carlo realisation of the mean-density
field, weighted so that the mean densities of the randoms and the data match.
Then the binned Landy--Szalay estimator for a separation bin $S$ is
\begin{equation}\label{eq:ls-slepian}
  \hat\xi(S)
  \;=\; \frac{\int \mathrm{d}\br_1 \mathrm{d}\br_2\;
              \theta(\br_1 - \br_2; S)\, N(\br_1) N(\br_2)}
             {\int \mathrm{d}\br_1 \mathrm{d}\br_2\;
              \theta(\br_1 - \br_2; S)\, R(\br_1) R(\br_2)} ,
\end{equation}
where $\theta(\cdot\,; S)$ is the bin indicator, equal to one when its
argument is a separation inside bin $S$ and zero otherwise.  Changing
variables to the separation $\bs = \br_2 - \br_1$ and
doing the integral over $\br_1$, which is a translation average, this
becomes
\begin{equation}\label{eq:ls-convolution}
  \hat\xi(S) \;=\;
  \frac{\int \mathrm{d}\bs\; \theta(\bs; S)\, [N \corr N](\bs)}
       {\int \mathrm{d}\bs\; \theta(\bs; S)\, [R \corr R](\bs)} ,
\end{equation}
where $\corr$ denotes correlation, the same shift-multiply-and-sum operation
as in Eq.~\eqref{eq:estimator}.

Read Eq.~\eqref{eq:ls-convolution} carefully, because it is the whole idea in
one line.  The Landy--Szalay estimator is a ratio of two correlations, each
subsequently averaged over a separation bin.  Both correlations can be
computed by fast Fourier transform.  Nothing has been approximated to get
here.

Our Eq.~\eqref{eq:ls-equals-ours} is the special case in which $R$ is
constant inside the mask, so that $[R\corr R](d)$ reduces to
$\bar\rho^2 n(d)$ and can be written down without transforming anything.
When the selection is not uniform, the recipe is unchanged except that the
denominator must be computed the same way as the numerator, by transforming
$R$.  The identity between the two methods, which is the subject of this note,
holds either way.

There is one important consequence of writing the estimator with $N$ and $R$
rather than with $\delta$, and it is easy to miss.  Because the padded zeros
enter the numerator and the denominator in the same way, they cancel.  If
instead one forms $\delta$, pads it, transforms it and stops there ---
dividing by a constant volume rather than by the matching $[R\corr R]$ ---
the result is \emph{misnormalised} by the padding
\citep{Slepian2016}.  What protects Eq.~\eqref{eq:estimator} from that fate is
precisely the lag-dependent $n(d)$ in its denominator.  This is the single
most common way to get an FFT correlation function wrong, and
\S\ref{sec:denominator} is therefore not an implementation detail.

\section{The pair count for a rectangular region}\label{sec:denominator}

Both methods need $n(d) = \sum_x \W_x \W_{x+d}$, the number of valid pairs,
so we evaluate it once here before splitting into the two methods.

In general $n(d)$ is the autocorrelation of the mask, and one can obtain it
exactly as one obtains the numerator, by transforming $\W$ and inverting
$|\hat\W_k|^2$.  For a rectangular region we can do better and write it
down.

In one dimension with a solid run of $N$ occupied cells and a real edge, the
pairs at lag $d$ are those with both $x$ and $x+d$ in $[0,N)$, so
$x$ ranges over $[0, N-d)$ and
\begin{equation}\label{eq:nd-nonperiodic}
  n(d) \;=\; N - d, \qquad 0 \le d < N .
\end{equation}
On a periodic axis every cell has a partner at every lag, so
\begin{equation}\label{eq:nd-periodic}
  n(d) \;=\; N \quad \text{for all } d .
\end{equation}
The linear decline of Eq.~\eqref{eq:nd-nonperiodic} is the familiar
triangular pair-count profile.  It says that large separations are sampled by
fewer pairs than small ones, purely because of geometry, and dividing by it
is what stops that geometric effect being read as a decline in clustering.

In three dimensions the crucial property is that a rectangular mask is
\emph{separable}: it factorises into one indicator function per axis,
$\W_{\bx} = \W^{(1)}_{x_1}\W^{(2)}_{x_2}\W^{(3)}_{x_3}$.  The
autocorrelation of a product of functions of independent variables is the
product of their autocorrelations, so
\begin{equation}\label{eq:nd-3d}
  n(\bd) \;=\; n^{(1)}(d_1)\, n^{(2)}(d_2)\, n^{(3)}(d_3),
\end{equation}
with each factor given by Eq.~\eqref{eq:nd-nonperiodic} or
Eq.~\eqref{eq:nd-periodic} according to that axis's boundary condition.  For a
region that is periodic in $z$ and cut in $x$ and $y$,
\begin{equation}\label{eq:nd-pencil}
  n(\bd) \;=\; (N_x - d_1)(N_y - d_2)\,N_z .
\end{equation}

This closed form is why the direct lag sum can divide by a simple product and
still be a correct Landy--Szalay estimator, and it is also why the two
methods can share a denominator exactly.  Both use Eq.~\eqref{eq:nd-3d};
neither computes $RR$ numerically.

It is worth being explicit about the limitation.  Equation~\eqref{eq:nd-3d}
holds because the region is a rectangle.  For an irregularly shaped mask ---
an ellipsoidal Lagrangian patch, a survey footprint, a zoom region with a
ragged boundary --- there is no product formula, and $n(\bd)$ must be
computed as the autocorrelation of $\W$, most cheaply by the same
transform machinery used for the numerator.  The identity between the two
methods survives unchanged; only the closed form for the denominator is lost.

\section{Direct evaluation: the lag sum}\label{sec:direct}

The first way to evaluate $s(d)$ is to do what Eq.~\eqref{eq:estimator}
says.  For a given separation $d$, loop over every cell $x$ such that both $x$
and $x+d$ lie inside the region, accumulate $\delta_x\delta_{x+d}$, and count
the terms.

If the region has $N_{\rm cell}$ cells and we want all separations out to a
maximum lag of $k_{\max}$ cells per axis, the cost is
$\mathcal{O}(k_{\max}^3\, N_{\rm cell})$, because there are of order
$k_{\max}^3$ separation vectors and each one costs a pass over the grid.  For
a $16\times16\times128$ pencil with $k_{\max}=25$ this is roughly
$5\times10^{8}$ multiply-adds, which is seconds.  For a large region with
large $k_{\max}$ it becomes the dominant cost of the whole analysis, and
\S\ref{sec:cost} makes that comparison quantitative.

The virtue of this method is that there is nothing to get wrong
conceptually.  It is the definition, transcribed.  That makes it the natural
reference implementation against which anything faster should be checked.

\section{Evaluation by fast Fourier transform}\label{sec:fft}

\subsection{The discrete Fourier transform and its orthogonality relation}

Work in one dimension for the moment; the three-dimensional case is the same
statement applied on each axis, and we return to it at the end of
\S\ref{sec:wk}.  Let $f_x$ be an array of $M$ numbers indexed by
$x = 0,1,\dots,M-1$.  Its discrete Fourier transform (DFT) is
\begin{equation}\label{eq:dft}
  \hat f_k \;=\; \sum_{x=0}^{M-1} f_x\, e^{-2\pi i k x / M},
  \qquad k = 0,1,\dots,M-1 ,
\end{equation}
and the inverse transform is
\begin{equation}\label{eq:idft}
  f_x \;=\; \frac{1}{M} \sum_{k=0}^{M-1} \hat f_k\, e^{+2\pi i k x / M} .
\end{equation}
The fast Fourier transform (FFT) is an algorithm for evaluating
Eq.~\eqref{eq:dft} in $\mathcal{O}(M\log M)$ operations instead of the
$\mathcal{O}(M^2)$ that the formula suggests \citep{CooleyTukey1965}.  It
computes the same quantity; it is not a different or approximate transform.
This point deserves emphasis because it is the source of much of the
suspicion the Fourier method attracts.  An FFT is a regrouping of the sum in
Eq.~\eqref{eq:dft} using the algebraic structure of the roots of unity, in
the same way that Horner's rule is a regrouping of a polynomial evaluation.

The one fact we need is the \emph{orthogonality relation}
\begin{equation}\label{eq:orthogonality}
  \frac{1}{M}\sum_{k=0}^{M-1} e^{2\pi i k m / M}
  \;=\;
  \begin{cases}
    1 & \text{if } m \equiv 0 \pmod M,\\
    0 & \text{otherwise}.
  \end{cases}
\end{equation}
The proof is a geometric series and takes one line.  Write
$z = e^{2\pi i m / M}$.  If $m \equiv 0 \pmod M$ then $z = 1$ and the sum is
$M$, giving $1$ after the prefactor.  Otherwise $z \neq 1$, and the finite
geometric series gives $\sum_{k=0}^{M-1} z^k = (z^M - 1)/(z - 1)$.  The
numerator vanishes because $z^M = e^{2\pi i m} = 1$ for integer $m$, while
the denominator does not, so the sum is zero.

Note what Eq.~\eqref{eq:orthogonality} is: a Kronecker delta, but one that
identifies indices \emph{modulo} $M$.  That ``modulo'' is the whole story of
\S\ref{sec:wraparound}.

\subsection{Wiener--Khinchin as a finite identity}\label{sec:wk}

Now compute the inverse transform of the squared modulus, which is what
``transform, square, transform back'' means.  Start from the definition of
$|\hat f_k|^2$ and expand it as a double sum:
\begin{equation}\label{eq:modsq}
  |\hat f_k|^2
  \;=\; \hat f_k\, \overline{\hat f_k}
  \;=\; \left(\sum_x f_x e^{-2\pi i k x / M}\right)
        \left(\sum_y f_y e^{+2\pi i k y / M}\right)
  \;=\; \sum_x \sum_y f_x f_y\, e^{-2\pi i k (x-y)/M} ,
\end{equation}
where the overbar denotes complex conjugation, and we used that $f$ is real
so that $\overline{f_y} = f_y$.  Apply the inverse transform,
Eq.~\eqref{eq:idft}, and exchange the order of the sums, which is allowed
because all sums are finite:
\begin{align}
  \frac{1}{M}\sum_k |\hat f_k|^2 e^{2\pi i k d / M}
  &= \sum_x \sum_y f_x f_y \;
     \underbrace{\frac{1}{M}\sum_k e^{2\pi i k\,(d - x + y)/M}}_{\text{use Eq.~\eqref{eq:orthogonality}}}
     \label{eq:wk-step1}\\[4pt]
  &= \sum_x \sum_y f_x f_y \;
     \big[\, x - y \equiv d \!\!\pmod M \,\big]
     \label{eq:wk-step2}\\[4pt]
  &= \sum_x f_x\, f_{(x-d) \bmod M} .
     \label{eq:wk-step3}
\end{align}
In Eq.~\eqref{eq:wk-step2} the square bracket is $1$ when the condition
holds and $0$ otherwise.  Going to Eq.~\eqref{eq:wk-step3} we used that
condition to do the $y$ sum, which has exactly one surviving term for each
$x$, namely $y = (x-d) \bmod M$.  Relabelling the dummy index turns this into
$\sum_x f_x f_{(x+d) \bmod M}$, which is the more familiar form.

This is the Wiener--Khinchin theorem in its finite, discrete form.  Read it
carefully and notice what it is not.  There is no limit, no integral, no
truncation, no assumption of stationarity, no ensemble average.  It is an
identity between two finite sums of products of the same $M$ numbers, valid
for any array $f$ whatsoever.  In exact arithmetic the two sides are the same
number.

The generalisation to three dimensions is immediate: apply
Eq.~\eqref{eq:orthogonality} once per axis, and the three Kronecker deltas
together fix the full separation vector.  Nothing about the argument changes.

So the answer to ``how can squaring the modulus not lose information?'' is
this.  Squaring \emph{does} lose the phases, and the phases are exactly what
you would need to reconstruct the field $f$ itself.  But we never asked for
$f$.  We asked for its autocorrelation, and the autocorrelation is precisely
the phase-independent content of the field.  Nothing that $\xi$ contains has
been thrown away, because $\xi$ and $P(k)$ are a Fourier pair by definition.

\subsection{Circular versus linear autocorrelation}\label{sec:wraparound}

The identity just derived gives the \emph{circular} autocorrelation: the
index arithmetic is modulo $M$, so cell $M-1$ is treated as a neighbour of
cell $0$.  Whether that is what we want depends on the region.

If the axis in question is genuinely periodic --- a full simulation box, or
the long axis of a pencil beam that spans the box --- then wrapping is
physically correct.  A pair straddling the boundary is a real pair at a real
separation, and it should be counted.

If the axis was cut, so the region has a real edge there, wrapping is wrong.
It manufactures pairs between cells at opposite ends of the region, which are
not close together at all.

Take $M = N = 4$ with $f = (a,b,c,d)$ and look at lag $1$.  The circular
autocorrelation, which is what an unpadded FFT returns, is
\begin{equation}\label{eq:circ-example}
  s_{\rm circ}(1) \;=\; ab + bc + cd + \underbrace{da}_{\text{wrap}} .
\end{equation}
The linear autocorrelation, which is what a direct lag sum over an array with
real edges returns, is
\begin{equation}\label{eq:lin-example}
  s_{\rm lin}(1) \;=\; ab + bc + cd .
\end{equation}
They differ by the single term $da$, which pairs the last cell with the
first.  In general,
\begin{equation}\label{eq:circ-vs-lin}
  s_{\rm circ}(d) \;=\; s_{\rm lin}(d)
  \;+\; \sum_{x=N-d}^{N-1} f_x\, f_{x+d-N} ,
\end{equation}
and the extra sum has $d$ terms.  For small lags on a large array this is a
small contamination, which is exactly why it is easy to miss; for lags
approaching $N$ it dominates.

\subsection{How much padding is enough}\label{sec:pad-howmuch}

Now embed the $N$-cell array in a longer array of length $M$, filling the
extra entries with zeros:
\begin{equation}\label{eq:padded-array}
  \tilde f_x \;=\;
  \begin{cases}
    f_x & 0 \le x < N,\\
    0   & N \le x < M .
  \end{cases}
\end{equation}
Apply the identity of \S\ref{sec:wk} to $\tilde f$.  The wrap terms are those
in which one index exceeds $N-1$, and every such term contains a factor
$\tilde f$ evaluated in the padded zone, which is zero.  So they vanish
identically --- not approximately, but term by term.

How much padding is enough? A wrap term at lag $d$ pairs index $x$ near the
top of the array with index $x + d - M$.  For this to reach back into the
occupied region $[0,N)$ we need $x + d - M \ge 0$ with $x \le N-1$, that is
$d \ge M - N + 1$.  So all lags with
\begin{equation}\label{eq:pad-condition}
  d \;\le\; M - N
\end{equation}
are free of wrap contamination.

Equation~\eqref{eq:pad-condition} is the statement worth remembering, and it
is worth stating in words: \emph{pad by the largest separation you intend to
measure, not by the size of the region}.  If the region is $N$ cells across
but you only want $\xi$ out to $d_{\max} = N/4$, then
$M = N + N/4$ suffices, and the transform is done on an array a quarter
larger rather than twice as large.

The stronger-sounding rule that one must always embed a non-periodic region
in a padded array of twice its size is, as \citet{Slepian2016} point out, a
common misunderstanding.  It is the worst case of
Eq.~\eqref{eq:pad-condition}, obtained by demanding every lag the region can
supply, $d_{\max} = N$.  It is a safe default and it is what
\texttt{measure\_xi\_fft\_mpi.c} does, since for a pencil only $16$ cells
across there is nothing to gain by being clever.  For a large survey volume in
which only separations far below the box size are of interest, the saving
from padding by $d_{\max}$ instead is real.

The correct general statement is the one \citet{Slepian2016} give: the
periodic embedding need only be large enough that the separation between any
point in the region and any point in its periodic image exceeds the largest
separation being measured.

Take the $M=4$ example above and pad to $M=8$, giving
$\tilde f = (a,b,c,d,0,0,0,0)$.  The circular autocorrelation of $\tilde f$ at
lag $1$ is $ab + bc + cd + d\cdot 0 + 0 + 0 + 0 + 0\cdot a = ab+bc+cd$, which
is $s_{\rm lin}(1)$.  The offending $da$ term has been replaced by two terms
that each contain a zero.

Figure~\ref{fig:toy} works this through numerically, together with a small
two-dimensional grid in which one axis is cut and the other is periodic.  The
grid is deliberately small enough that any single entry can be recomputed by
hand from the definition, which is the only way to be sure that agreement
between two programs is not agreement between two copies of the same mistake.

\begin{figure}[htbp]
  \centering
  \includegraphics[width=\textwidth]{xi_toy.pdf}
  \caption{An $8\times16$ toy grid with the $x$ axis cut (zero-padded to
    $16$) and the $y$ axis periodic (not padded).  Left and centre: the
    correlation function evaluated at every lag by the direct sum and by the
    padded FFT.  The two panels are the same picture.  Right: the base-10
    logarithm of their absolute difference.  It sits between $10^{-16}$ and
    $10^{-14}$ everywhere, so the two evaluations differ only in the last few
    digits a double carries --- the residue \S\ref{sec:roundoff} accounts for.
    Produced by \texttt{figures/xi\_estimators/plot\_xi\_toy.py}, which also
    prints the four-cell one-dimensional example of
    Eqs.~\eqref{eq:circ-example}--\eqref{eq:lin-example} with its wrap-around
    term written out.}
  \label{fig:toy}
\end{figure}

\subsection{Mixed boundary conditions: the pencil beam}\label{sec:mixed}

Nothing in \S\ref{sec:wraparound} or \S\ref{sec:denominator} requires all
three axes to be treated alike, and a pencil beam is precisely the case where
they are not.  The region spans the full periodic box along one axis and is
cut on the other two.

The recipe follows directly.  For each axis independently:
\begin{itemize}
  \item \textbf{Periodic axis.} Do not pad.  The circular autocorrelation that
        the unpadded FFT returns is the physically correct one, because
        wrapping pairs are real.  Use $n = N$ on that axis.
  \item \textbf{Non-periodic axis.} Pad by at least $d_{\max}$, and to $2N$
        if every available lag is wanted.  The zeros suppress the wrap terms,
        so the FFT returns the linear autocorrelation.  Use $n = N - d$ on
        that axis.
\end{itemize}
The padded array for a $16\times16\times128$ pencil with $z$ periodic is
therefore $32\times32\times128$, not $32\times32\times256$.  Not padding the
long axis is a saving as well as a physical statement: the wrap along $z$
carries genuine information about pairs separated by more than half the
region, which is what lets a pencil beam reach separations far beyond half
its transverse width.

The direct lag sum makes the same distinction by choosing, per axis, whether
its index arithmetic wraps.  Because both methods make the choice the same way
and use the same $n(\bd)$, they agree axis by axis.
Figure~\ref{fig:pencil} shows the outcome at a realistic size.

\begin{figure}[htbp]
  \centering
  \includegraphics[width=\textwidth]{xi_identity.pdf}
  \caption{A $16\times16\times128$ cell pencil
    ($125\times125\times1000\ \mathrm{Mpc}/h$) cut from a $128^3$ periodic
    box, with $x$ and $y$ non-periodic and $z$ periodic.  \emph{Left:} the
    power spectrum.  The full box recovers the input spectrum, as it must.
    The pencil's $|\mathrm{FFT}(\delta_{\rm pad})|^2$ --- the intermediate the
    padded-FFT estimator forms on the way to $\xi$ --- is the true $P(k)$
    convolved with the mask window, so it is suppressed and scattered at low
    $k$ and converges to the input only well inside the region scale.
    \emph{Right:} the correlation function from both estimators.  The circles
    lie on the line at every point, and the lower panel shows the residual
    relative difference.  The departure of both from linear
    theory beyond $r \sim 80\ \mathrm{Mpc}/h$ is cosmic variance in a single
    pencil together with the integral constraint of
    \S\ref{sec:integral-constraint}, and it affects the two estimators
    identically.  Produced by
    \texttt{figures/xi\_estimators/plot\_xi\_identity.py}.}
  \label{fig:pencil}
\end{figure}

\section{Rounding error}\label{sec:roundoff}

If the two computations are the same expression, why do they not agree
bit-for-bit?

Because floating-point addition is not associative.  Both methods add up the
same set of products, but in different orders, and each ordering commits a
different sequence of rounding errors.  The exact answer is the same; the two
computed answers straddle it.

The two orderings have different error growth.  Write $\epsilon$ for the
machine epsilon, the spacing between representable numbers near unity, which
is $2.2\times10^{-16}$ in double precision.  The direct sum at a given lag adds
$n(\bd)$ products, one per valid pair.  Summing them sequentially
accumulates a relative error bounded by roughly $n(\bd)\,\epsilon$, and
in practice, with errors of random sign, of order
$\sqrt{n(\bd)}\,\epsilon$.  An FFT
of length $M$ performs $\log_2 M$ butterfly stages, and its error grows like
$\sqrt{\log_2 M}\,\epsilon$ \citep{Frigo2005}.  For large arrays the FFT is
therefore the \emph{better}-conditioned of the two, which is worth knowing:
the direct sum is the conceptual reference, not the numerically superior one.

This also makes the comparison a useful diagnostic.  A difference far above
the rounding level means one of the two programs is not evaluating the
estimator that was derived above, and the cause is essentially always one of
a short list: insufficient padding on a cut axis
(\S\ref{sec:pad-howmuch}); a denominator that is a constant volume rather
than the lag-dependent $n(\bd)$ (\S\ref{sec:denominator},
\S\ref{sec:ls-general}); a mass-assignment window deconvolved in one path but
not the other; or two different choices of mean density
(\S\ref{sec:integral-constraint}).

\section{Computational cost}\label{sec:cost}

The identity of \S\S\ref{sec:direct}--\ref{sec:fft} says the two methods give
the same answer.  It says nothing about which to run, and that is decided
entirely by cost.

Direct pair counting scales as the square of the number of objects within the
maximum separation of interest.  For a catalogue of $N$ objects with number
density $n$, counting all pairs out to a radius enclosing volume $V_{\max}$
costs $\mathcal{O}(N n V_{\max})$, which for a survey is
$\mathcal{O}\big((nV)^2\big)$ in the worst case \citep{Slepian2016}.  Trees
and grids reduce the constant by avoiding pairs far beyond the scale of
interest, but they do not change the scaling, and they help least when the
binning is coarse compared with the interparticle spacing --- which is the
usual situation.

The Fourier method costs $\mathcal{O}(N_g \log N_g)$ in the number of grid
cells $N_g$, plus a term linear in the number of objects for the mass
assignment itself.  The two costs are therefore controlled by different
things: pair counting by how many objects and how large a separation, the
Fourier method by how fine a grid.

Three regimes follow, and they are worth stating because they say when the
work of implementing the Fourier method pays for itself
\citep{Slepian2016}:
\begin{itemize}
  \item \textbf{Large separations.} Pair counting cost grows with
        $V_{\max}$; FFT cost does not depend on the separation range at all.
        The larger the scales of interest, the more the Fourier method wins.
  \item \textbf{High number density.} Pair counting cost grows quadratically
        with density; the FFT cost is fixed by the grid, so it is unchanged
        when the same volume holds ten times as many objects.
  \item \textbf{Coarse grids.} If the science can tolerate a coarse grid ---
        which is a real constraint, see \S\ref{sec:gives-up} --- $N_g$ is
        small and the FFT is very cheap.
\end{itemize}
Conversely, if you want small separations in a sparse catalogue at high
spatial resolution, pair counting may well be the cheaper option, and it
avoids the artefacts of the next section entirely.

There is a second, more mundane reason the scaling matters.  Modern analyses
compute covariance matrices from hundreds or thousands of mock catalogues,
and that repetition, not the single measurement on real data, dominates the
total cost.  A method that is somewhat less accurate but far faster is often
the right trade for the mocks even when the real data are analysed with an
explicit pair count \citep{Slepian2016}.

\section{Errors introduced by the grid}\label{sec:gives-up}

The identity proved above is exact \emph{given the gridded field}.  It says
nothing about the step before that, and it is there that the Fourier method
concedes something real.  Three concessions are worth knowing
\citep{Slepian2016}.

\textbf{Gridding displaces the objects.} Assigning particles to cells moves
each one to a nearby cell centre, or spreads it over neighbouring cells.  The
measured correlation function is therefore a smoothed version of the truth,
with the smoothing set by the assignment scheme and the cell size
\citep{Jing2005}.  Pair counting uses the exact positions and has no such
error.  The size of the effect is governed by the ratio of the grid spacing to
the width of the radial bins: if the bins are much wider than a cell, the
smoothing is buried inside the bin and does not matter.

A worked figure makes this concrete.  A baryon acoustic oscillation analysis
using separation bins of $8\,h^{-1}$~Mpc wants a grid comfortably finer than
that.  \citet{Slepian2016} suggest tuning the bins so that they have
negligible support beyond $k \approx 0.4\,h\,\mathrm{Mpc}^{-1}$ and then
choosing a grid scale of $3$--$4\,h^{-1}$~Mpc, which places the Nyquist
frequency safely above the smoothing scale.

\textbf{A Cartesian grid does not align with spherical bins.} Our lags live
on a cubic lattice, but $\xi(r)$ is a function of $|r|$, so each radial bin
is assembled from whichever lattice vectors happen to fall inside it.  At
small $r$ there are few such vectors and the binning is visibly ragged; this
is why the innermost bins in Fig.~\ref{fig:pencil} are built from as few as
one or three lags.  Pair counting in spherical coordinates has no analogue of
this.  One partial remedy is to use radial bins with smooth edges rather than
top hats, which both reduces the artefacts and is numerically better behaved
when the result is to be compared against a theoretical spectrum through a
spherical Bessel transform.

\textbf{Normalisation is only right if the denominator matches.} This is the
point already made in \S\ref{sec:ls-general}, and it belongs on any list of
what can go wrong: transforming a padded $\delta$ field and dividing by a
constant gives a misnormalised answer, because the padding was never
accounted for.  Dividing instead by the matching pair count, whether the
analytic $n(\bd)$ of Eq.~\eqref{eq:nd-3d} or a transformed $R$ field,
is what makes the padding cancel.

The reassuring statement to hold onto is that all three concessions are
controlled by the grid, and all vanish as the grid is refined: the FFT
estimate converges to the pair-counting answer as $N_g$ grows
\citep{Slepian2016}.  They are a resolution choice, not a bias that cannot be
removed.

\section{The mean density and the integral constraint}\label{sec:integral-constraint}

One input deserves separate attention because it is easy to get subtly wrong
and because it does not announce itself as an error.

The field $\delta$ depends on the mean density $\bar\rho$ chosen in
Eq.~\eqref{eq:delta-def}, and $\hat\xi$ depends on $\delta$ quadratically.
The mean must be taken over the \emph{region}, that is over cells with
$\W_x = 1$, and not over the padded array.  The padded array contains the
zeros we inserted, and its mean is smaller by the volume ratio
$V_{\rm region}/V_{\rm padded}$.  Using it would rescale the field and shift
$\xi$ by an amount that is not a rounding error.

With the mean taken over the region, $\sum_x \W_x\delta_x = 0$ by
construction, so the padded field has zero total sum and its $k=0$ Fourier
mode vanishes.  This is the discrete statement of the \emph{integral
constraint}: because the mean was estimated from the same data, the measured
$\xi$ is forced to integrate to zero over the region, and it is therefore
biased low relative to the true $\xi$ by an amount of order the variance of
the region's mean density \citep{Peebles1980}.  That bias is a property of
the estimator, shared identically by both methods, and it does not affect
their agreement with each other.  It does affect comparison with theory, which
is why a measured $\xi$ from a small region should not be expected to sit on
a linear-theory curve at separations approaching the region size --- the
effect visible at large $r$ in Fig.~\ref{fig:pencil}.

Because the integral constraint is a property of the Landy--Szalay estimator
rather than of how its sums are evaluated, it is unchanged by moving to the
Fourier method \citep{Slepian2016}.  Whatever was true of the pair-counted
monopole remains true of the transformed one.

\section{What a comparison of the two methods tests}\label{sec:validates}

Suppose you run both estimators on the same data and they agree.  What have
you learned?

You have learned that the fast implementation faithfully reproduces the
reference implementation of the correlation step: the transform, the padding,
the lag indexing, the denominator lookup, and the binning.  Given how many
sign and index conventions are involved, this is a genuinely strong check,
and it is the right check to run before trusting the fast path.

You have \emph{not} learned that the estimator is correct.  If the two
programs share their upstream code --- the same mass assignment, the same
mean-density normalisation, the same buffer-cell handling, the same $r$-bin
edges --- then any error in that shared code appears identically in both
outputs and cancels exactly out of the comparison.  A comparison can only test
what differs between the two things being compared.

Nor does it test anything in \S\ref{sec:gives-up}.  The gridding artefacts are
present in both methods by construction, because both start from the same
gridded field.  Detecting those requires either refining the grid and watching
the answer move, or comparing against a genuine pair count on the ungridded
positions.

Validating the shared part therefore requires a different test with an
external reference.  Useful options are: run on a field with an analytically
known correlation function and compare; run on a full periodic box, where the
answer can be checked against an independent full-box estimator; refine the
grid and confirm convergence; or compare against linear theory at separations
well inside the region, where the integral constraint of
\S\ref{sec:integral-constraint} is small.

\section{Summary}

\begin{enumerate}
  \item A grid of $\delta$ with a lag-dependent pair-count denominator is the
        Landy--Szalay estimator in the infinite-random limit
        (Eq.~\ref{eq:ls-equals-ours}).  The $RR$ term has been evaluated in
        closed form, not neglected.  The general form for non-uniform
        selection is Eq.~\eqref{eq:ls-convolution}: a ratio of two
        correlations, both computable by FFT \citep{Slepian2016}.
  \item Transforming, squaring the modulus, and transforming back returns the
        autocorrelation exactly (Eq.~\ref{eq:wk-step3}).  This is a finite
        algebraic identity following from DFT orthogonality, not a numerical
        approximation.  Squaring discards the phases, which are needed to
        reconstruct the field but not its correlation function.
  \item The FFT returns the \emph{circular} autocorrelation.  Zero-padding a
        cut axis removes the wrap terms exactly, because each such term
        contains a factor drawn from the padded zone.  The requirement is
        $d \le M-N$ (Eq.~\ref{eq:pad-condition}): pad by the largest
        separation wanted, not necessarily by the size of the region.  A
        genuinely periodic axis is left unpadded, because there the wrap is
        physical.
  \item For a rectangular region the pair count factorises per axis into
        $N$ or $N-d$ (Eq.~\ref{eq:nd-3d}), so both methods can divide by the
        same analytic denominator.  An irregular mask loses the closed form
        but not the identity.  Dividing by a constant volume instead is the
        standard way to get a misnormalised answer.
  \item The two methods therefore compute the same expression and differ only
        in summation order, so they agree to rounding.  A discrepancy above
        that level points to padding, denominator, window, or mean-density
        mismatch.
  \item Which method to \emph{run} is a question of cost, not correctness.
        The Fourier method wins for large separations, high number density and
        coarse grids; pair counting wins for small separations in sparse
        catalogues, and avoids gridding artefacts entirely.
  \item What the Fourier method concedes is all upstream of the identity:
        gridding smooths the measured $\xi$, a Cartesian lattice bins
        awkwardly into spherical shells, and the normalisation is only
        correct if the denominator matches the numerator.  All three are
        controlled by the grid and vanish as it is refined.
  \item Agreement between the two methods validates the correlation
        machinery and says nothing about code shared by both paths, nor
        about any of the concessions above.
\end{enumerate}

\bibliography{references}

\end{document}
